\section{Introduction}

Personalized text-to-image generation has progressed rapidly in recent years, especially in the single-subject regime where a model must preserve one reference identity while following a textual prompt. However, practical deployment increasingly requires multiple distinct reference identities to co-exist in one image, often with explicit interactions, physical contact, asymmetric role assignments, or object bindings. This multi-subject setting is substantially harder than ordinary composition because the model must preserve identity boundaries, local visual traits, and subject-to-action bindings at the same time.

The evaluation layer has not kept pace with this difficulty. In many current pipelines, quality is still measured with global similarity metrics or general preference models. These tools are useful for coarse semantic plausibility, but they are poorly matched to the real errors that dominate multi-subject personalization. A generated image can mention the right scene, contain approximately the right colors, and still be unusable because one subject is missing, two identities collapse into generic clones, clothing attributes bleed across people, or a requested interaction is assigned to the wrong subject. In such cases, a useful evaluator must preserve ranking separability even when the number of requested subjects becomes large.

This paper is motivated by a simple empirical observation: once the subject count increases, existing metrics often lose the ability to separate obviously better generations from obviously worse ones. The problem is therefore not only that current metrics are imperfect, but that they become structurally uninformative in the regime that matters most for scalable personalization. If a benchmark cannot distinguish ``missing one subject'' from ``missing most subjects,'' then it cannot support real model development or credible ablation.

We formalize this failure mode through two coupled contributions. First, we introduce \textbf{PrismBench}, a benchmark specifically designed for reference-conditioned multi-subject generation. Each task consists of a prompt, a set of reference subjects, and a paired comparison between two generated images produced under the same prompt and references. The current benchmark construction uses a 60k-pair silver pipeline built from Nano and Mosaic generations labeled by Gemini 2.5 Flash and Gemini 3.1 Flash Lite, together with a 4k-pair human evaluation set spanning Nano, Mosaic, Flux, Seedream 4.5, GPT-Image-1.5, and GLM. Second, we propose \textbf{LENS}, a learned evaluator that takes the references, prompt, and generated image pair as input, then predicts both pairwise preference and three interpretable image-level diagnostic dimensions: \emph{Existence}, \emph{Appearance}, and \emph{Interaction}.

The design choices in PrismBench and LENS are driven by a practical principle: preserve maximal ranking signal while minimizing annotation burden. Instead of requiring annotators to label the exact number of missing subjects or to score each subject with fine-grained severity levels, we adopt binary image-level diagnosis and a simple A/B preference decision. This makes large-scale annotation feasible, improves consistency, and still enables the score head to learn a continuous severity scale from pairwise comparisons. The silver pipeline is intentionally conservative: preference conflicts are dropped, while diagnostic disagreement is retained as averaged soft supervision.

The expected outcome is not merely another scalar metric, but a benchmark-and-metric stack that can answer three questions at once: which image is better, why it is better, and whether the metric remains discriminative when subject count and interaction difficulty increase. The paper is therefore centered on a broader thesis: \emph{multi-subject personalized generation requires a new evaluation protocol, not just a stronger backbone}. The experimental program then becomes straightforward: train LENS on the large silver set, stress-test it across backbone scales and update modes, and compare it against standard baselines on a broader human evaluation set collected from multiple generator families.

The remainder of the paper is organized as follows. Section~2 reviews the limitations of existing evaluation paradigms. Section~3 defines the problem setup and diagnostic taxonomy. Section~4 introduces PrismBench, including data collection and annotation protocols. Section~5 presents the LENS evaluator. Section~6 describes the training pipeline. Section~7 outlines the planned experiments and success criteria, and Section~8 concludes with discussion and limitations.
